\chapter{ระบบ AIoT และการประมวลผลปัญญาประดิษฐ์ฝังตัวบนขอบ}

การเชื่อมต่อโลกเสมือนจริงเข้ากับเซนเซอร์ภาคสนาม เป็นแกนกลางของการเกษตรแม่นยำ (Precision Agriculture) และฟิสิกส์ประยุกต์ โดยใช้โมเดลปัญญาประดิษฐ์ประมวลผล ณ จุดกำเนิดข้อมูล (Edge AI)

\section{การสื่อสารฮาร์ดแวร์เซนเซอร์วัดค่าดินผ่าน RS485 Modbus RTU}

หัววัดคุณสมบัติดินแบบแท่งโลหะหลายพารามิเตอร์ (Multi-parameter Soil Probe) ส่งข้อมูลผ่านโปรโตคอลอุตสาหกรรม Modbus RTU ด้วยความเร็ว 9600 bps โดยมีโครงสร้างเฟรมคำสั่งมาตรฐาน (Request Frame) ดังนี้:

\begin{itemize}
    \item \textbf{Slave ID (1 Byte):} รหัสประจำตัวอุปกรณ์ เช่น \texttt{0x01}
    \item \textbf{Function Code (1 Byte):} ฟังก์ชันอ่านค่ารีจิสเตอร์ \texttt{0x03}
    \item \textbf{Starting Address (2 Bytes):} ตำแหน่งแอดเดรสเริ่มต้น เช่น \texttt{0x0000}
    \item \textbf{Quantity of Registers (2 Bytes):} จำนวนรีจิสเตอร์ที่ต้องการอ่าน
    \item \textbf{CRC16 Checksum (2 Bytes):} รหัสตรวจสอบความถูกต้องของข้อมูล
\end{itemize}

\section{การฝึกสอนโมเดล Ordinary Least Squares บน Edge Device}

ในการปรับเทียบความถูกต้องของค่าความเป็นกรด-ด่าง (pH) และความชื้นในดิน อัลกอริทึมการถดถอยเชิงเส้นพหุคูณคำนวณผ่านเวกเตอร์เมทริกซ์ในหน่วยความจำของบอร์ดประมวลผล:

\begin{equation}
\hat{\beta} = (\mathbf{X}^T \mathbf{X})^{-1} \mathbf{X}^T \mathbf{y}
\end{equation}

โดยที่ $\mathbf{X}$ คือเมทริกซ์ของสัญญาณไฟฟ้าและอุณหภูมิ และ $\mathbf{y}$ คือค่ามาตรฐานที่วัดได้จากห้องปฏิบัติการเคมีมาตรฐาน ผลลัพธ์ที่ได้มีความแม่นยำสูงและตอบสนองได้แบบเรียลไทม์โดยไม่ต้องพึ่งพาเซิร์ฟเวอร์คลาวด์ภายนอก
